% ── chemistry-handout-preamble.tex ── 化学竞赛讲义 LaTeX 导言区（v6）
\documentclass[a4paper,11pt]{ctexart}

% ── 字体优化 ──
% Windows 标准字体：SimSun 宋体正文，SimHei 黑体标题
% 无需检测 Noto，直接使用 Windows 自带字体（稳定 + 快速）
\setCJKmainfont[AutoFakeBold=true,BoldFont=SimHei]{SimSun}
\setCJKsansfont{SimHei}

% 化学式宏（mhchem）
\usepackage[version=4]{mhchem}

% ── 数学 ──
\usepackage{amsmath,amssymb}

% ── 图片 ──
\usepackage{graphicx}

% ── 表格（跨页+自动列宽）──
\usepackage{booktabs}
\usepackage{xltabular}
\usepackage{array}
\newcolumntype{L}{>{\raggedright\arraybackslash}X}
\newcolumntype{C}{>{\centering\arraybackslash}X}
\newcolumntype{R}{>{\raggedleft\arraybackslash}X}

% ── 浮动体（[H] 强制定位）──
\usepackage{float}

% ── 页面 ──
\usepackage[includeheadfoot,margin=2.2cm,headheight=14pt]{geometry}
\setlength{\textwidth}{16.6cm}

% ── 超链接(无红框) ──
\usepackage[colorlinks=true,linkcolor=black,citecolor=black,urlcolor=black]{hyperref}

% ── 颜色 ──
\usepackage[dvipsnames,svgnames]{xcolor}
\definecolor{chapterblue}{RGB}{23,50,77}
\definecolor{insightbg}{RGB}{235,240,250}
\definecolor{warningbg}{RGB}{255,245,235}
\definecolor{tipbg}{RGB}{235,250,235}

% ── 页眉页脚 ──
\usepackage{fancyhdr}
\pagestyle{fancy}
\fancyhf{}
\fancyhead[L]{\small\color{chapterblue}\leftmark}
\fancyhead[R]{\small\color{chapterblue}\thepage}
\renewcommand{\headrule}{{\color{chapterblue}\hrule height 0.8pt}}

% ── 章节编号与样式 ──
\setcounter{secnumdepth}{1}
\ctexset{
  section = {
    format = \Large\bfseries\heiti\color{chapterblue},
    aftername = \hspace{0.5em},
    beforeskip = 2ex,
    afterskip = 1ex,
  },
  subsection = {
    format = \large\bfseries\heiti,
    beforeskip = 1.5ex,
    afterskip = 0.5ex,
  },
}

% ── 教学提示框（浅底色版）──
\newenvironment{insightbox}{%
  \par\smallskip
  \noindent{\small\bfseries\color{RoyalBlue}🧠 认知冲突}\par
  \leavevmode\colorbox{insightbg}{\parbox{\dimexpr\textwidth-2\fboxsep\relax}{%
  \ignorespaces}}%
}{\par\smallskip}

\newenvironment{warningbox}{%
  \par\smallskip
  \noindent{\small\bfseries\color{BrickRed}⚠️ 易错点}\par
  \leavevmode\colorbox{warningbg}{\parbox{\dimexpr\textwidth-2\fboxsep\relax}{%
  \ignorespaces}}%
}{\par\smallskip}

\newenvironment{tipbox}{%
  \par\smallskip
  \noindent{\small\bfseries\color{ForestGreen}💡 理解提示}\par
  \leavevmode\colorbox{tipbg}{\parbox{\dimexpr\textwidth-2\fboxsep\relax}{%
  \ignorespaces}}%
}{\par\smallskip}

% ── 列表（答案区加间距）──
\usepackage{enumitem}
\setlist{nosep}
% 专门用于参考答案的列表（每项之间空一行）
\newlist{answerlist}{enumerate}{1}
\setlist[answerlist]{label=\arabic*., leftmargin=2em, itemsep=0.8em}

% ── 防止内容溢出 ──
\usepackage{microtype}
\setlength{\emergencystretch}{2em}
\setlength{\hfuzz}{2pt}
\setlength{\overfullrule}{0pt}
\sloppy

% ── pandoc 兼容 ──
\providecommand{\tightlist}{\setlength{\itemsep}{0pt}\setlength{\parskip}{0pt}}
\newenvironment{Shaded}{}{}
\newenvironment{Highlighting}{}{}
\newcommand{\NormalTok}[1]{#1}
\newcommand{\CommentTok}[1]{#1}
\newcommand{\KeywordTok}[1]{#1}
\newcommand{\StringTok}[1]{#1}
\newcommand{\FunctionTok}[1]{#1}
\newcommand{\DataTypeTok}[1]{#1}
\newcommand{\NumberTok}[1]{#1}
\newcommand{\ControlFlowTok}[1]{#1}
\newcommand{\OperatorTok}[1]{#1}
\newcommand{\BaseNTok}[1]{#1}
\newcounter{none}

% ── 图片尺寸限制 ──
\AtBeginDocument{\setkeys{Gin}{width=\textwidth,keepaspectratio}}

% ── 防止图片跨节 ──
\usepackage{placeins}

% ── 自定义命令 ──
% 练习题标题（去掉括号）
% 例题标题（去掉中括号）
% 这些需要在convert.py预处理中做正则替换

\graphicspath{{C:/Temp/chem_media/}}
\title{晶体学基础}
\date{}
\setlength{\tabcolsep}{3.5pt}
\small
\begin{document}
\maketitle
\tableofcontents
\newpage
\section{第六讲：晶体学基础}\label{ux7b2cux516dux8bb2ux6676ux4f53ux5b66ux57faux7840}

\textbf{本讲定位}：系统建立晶体学的概念语言------从宏观晶体特征到微观周期性排列，掌握点阵、晶胞、晶系与对称性的数学描述。本讲是后续金属晶体、离子晶体、原子晶体等各类晶体结构学习的语言基础。

\begin{center}\rule{0.5\linewidth}{0.5pt}\end{center}

\subsection{学习目标}\label{ux5b66ux4e60ux76eeux6807}

\begin{itemize}
\tightlist
\item
  ✅ 能说出晶体的四个基本性质（自范性、各向异性、对称性、均一性）
\item
  ✅ 能理解点阵与结构基元的关系（晶体结构 = 点阵 + 结构基元），能辨认简单晶体的结构基元
\item
  ✅ 能根据轴长和轴角特征判断七大晶系
\item
  ✅ 能识别14种布拉维格子（按晶系+带心方式）的基本特征
\item
  ✅ 能计算晶胞中的原子数（顶角 \(1/8\)、棱心 \(1/4\)、面心 \(1/2\)、体内 \(1\)）
\item
  ✅ 能区分素晶胞与复晶胞（P、I、F、C），并能用平移操作检验晶胞类型
\item
  ✅ 能用分数坐标描述晶胞中原子的位置
\item
  ✅ 能识别分子/晶体中的基本对称元素（旋转轴 \(C_n\)、镜面 \(\sigma\)、对称中心 \(i\)）
\end{itemize}

\begin{center}\rule{0.5\linewidth}{0.5pt}\end{center}

\subsection{一、晶体与晶体的基本性质}\label{ux4e00ux6676ux4f53ux4e0eux6676ux4f53ux7684ux57faux672cux6027ux8d28}

\subsubsection{1.1 晶体的定义}\label{ux6676ux4f53ux7684ux5b9aux4e49}

晶体是原子、离子或分子按照一定的周期性在空间排列，形成具有一定规则几何外形的固体。与非晶态物质（玻璃、树脂）相比，晶体具有长程有序的微观结构。

\begin{quote}
💡 \textbf{深度辨析}：「液晶 (Liquid Crystal)」------介于晶体与液体之间的中间态。液晶分子具有取向有序（像晶体）但位置无序（像液体）。液晶的发现挑战了''晶体 vs 非晶体''的二元分类，也因此获得了1991年诺贝尔物理学奖（de Gennes）。液晶显示器(LCD)正是利用液晶的电光效应工作的。
\end{quote}

\subsubsection{1.2 晶体的四个基本性质}\label{ux6676ux4f53ux7684ux56dbux4e2aux57faux672cux6027ux8d28}

{\def\LTcaptype{none} % do not increment counter
\begin{longtable}[]{@{}
  >{\raggedright\arraybackslash}p{(\linewidth - 4\tabcolsep) * \real{0.3333}}
  >{\raggedright\arraybackslash}p{(\linewidth - 4\tabcolsep) * \real{0.3333}}
  >{\raggedright\arraybackslash}p{(\linewidth - 4\tabcolsep) * \real{0.3333}}@{}}
\toprule\noalign{}
\begin{minipage}[b]{\linewidth}\raggedright
性质
\end{minipage} & \begin{minipage}[b]{\linewidth}\raggedright
定义
\end{minipage} & \begin{minipage}[b]{\linewidth}\raggedright
实例说明
\end{minipage} \\
\midrule\noalign{}
\endhead
\bottomrule\noalign{}
\endlastfoot
\textbf{自范性} & 在适宜条件下，晶体能自发地呈现封闭的、规则的凸多面体外形 & 氯化钠晶体呈立方体，石英晶体呈六方柱 \\
\textbf{各向异性} & 沿晶体的不同方向，物理性质（折射率、热导率、解理性等）不同 & 云母沿层状方向易剥离，垂直方向难；石墨沿层方向导电，垂直方向不导电 \\
\end{longtable}
}

\begin{quote}
🧠 \textbf{认知冲突}：「玻璃是固体，为什么不是晶体？」------因为玻璃中的原子/分子排列是\textbf{短程有序但长程无序}（像冻结的液体）。它没有固定的熔点，只是逐渐软化。这揭示了晶体的本质特征不是''硬''或''固体''，而是\textbf{周期性长程有序}。
\end{quote}

\textbf{对称性} \textbar{} 晶体的理想外形和内部结构具有对称性 \textbar{} 立方晶体绕 \(C_4\) 轴旋转 \(90^\circ\) 后与自身重合 \textbar{}\\
\textbf{均一性} \textbar{} 同一晶体的各个部分宏观性质相同 \textbar{} 晶体内任意取两点，化学组成和物理性质一致 \textbar{}

\textbf{注意}：均一性并不与各向异性矛盾------均一性指同一方向上的性质各处相同，各向异性指不同方向上的性质可能不同。

\subsubsection{1.3 晶体与非晶体的对比}\label{ux6676ux4f53ux4e0eux975eux6676ux4f53ux7684ux5bf9ux6bd4}

{\def\LTcaptype{none} % do not increment counter
\begin{longtable}[]{@{}lll@{}}
\toprule\noalign{}
特征 & 晶体 & 非晶体（玻璃态） \\
\midrule\noalign{}
\endhead
\bottomrule\noalign{}
\endlastfoot
微观结构 & 长程有序（周期性排列） & 短程有序、长程无序 \\
宏观外形 & 规则几何多面体 & 无规则外形 \\
熔点 & 固定（有明确的熔点） & 不固定（软化温度范围） \\
各向异性 & 显著 & 各向同性 \\
X射线衍射 & 有尖锐衍射峰 & 无尖锐衍射峰（弥散散射） \\
\end{longtable}
}

\begin{figure}[H]\centering
\includegraphics[width=0.9\textwidth,keepaspectratio]{13-1-quartz-crystal-faces-angle-constancy.png}
\caption{图 石英晶体的不同外形与晶面——虽然外形各异，但对应晶面间的夹角始终相等（晶面夹角守恒定律），这是晶体自范性的宏观体现（普化原理第4版 图13.1）}
\end{figure}

\begin{center}\rule{0.5\linewidth}{0.5pt}\end{center}

\subsection{二、点阵理论}\label{ux4e8cux70b9ux9635ux7406ux8bba}

\subsubsection{2.1 结构基元与点阵}\label{ux7ed3ux6784ux57faux5143ux4e0eux70b9ux9635}

晶体内部的质点在三维空间作周期性的有规则排列。为研究这种周期性，将晶体抽象为 \textbf{点阵（lattice）}。

\textbf{结构基元}：晶体中周期性变化的具体内容。各个结构基元必须满足四个条件：化学组成相同、空间结构相同、排列取向相同、周围环境相同。

\textbf{示例}：在伸展的聚乙烯链 \(-\mathrm{CH}_2-\mathrm{CH}_2-\) 中，两个相邻的 \(-\mathrm{CH}_2-\) 虽组成和结构相同，但取向不同、周围环境不同，因此一个结构基元应为 \(-\mathrm{CH}_2-\mathrm{CH}_2-\)，而非单个 \(-\mathrm{CH}_2-\)。

\begin{quote}
🧠 \textbf{认知冲突}：「为什么聚乙烯的结构基元不是单个CH\textsubscript{2}？」------单个CH\textsubscript{2}虽然组成和结构完全相同，但相邻两个CH\textsubscript{2}的C---C键取向相反（一个''正放''、一个''倒放''），周围环境也不相同（一个左边是CH\textsubscript{2}、右边是CH\textsubscript{2}，但取向不同）。所以必须取两个CH\textsubscript{2}作为一个结构基元。
\end{quote}

确定了结构基元后，用一个抽象的几何点（点阵点）代表每个结构基元。

\begin{quote}
📝 \textbf{练一练：结构基元识别}
\textbar{} 物质 \textbar{} 结构基元是？ \textbar{} 判断理由 \textbar{}
\textbar:---\textbar:---\textbar:---\textbar{}
\textbar{} NaCl晶体 \textbar{} \_\_\_\_\_\_ \textbar{} \_\_\_\_\_\_ \textbar{}
\textbar{} 干冰（CO\textsubscript{2}(s)） \textbar{} \_\_\_\_\_\_ \textbar{} \_\_\_\_\_\_ \textbar{}
\textbar{} 金属Cu \textbar{} \_\_\_\_\_\_ \textbar{} \_\_\_\_\_\_ \textbar{}
\textbar{} 伸展的聚乙烯链 \textbar{} \_\_\_\_\_\_ \textbar{} \_\_\_\_\_\_ \textbar{}
\textbar{} CsCl晶体 \textbar{} \_\_\_\_\_\_ \textbar{} \_\_\_\_\_\_ \textbar{}
\end{quote}

从晶体中无数结构单元中抽象出的几何点集合构成 \textbf{点阵（lattice）}。

\textbf{核心关系}：

\[
\text{晶体结构} = \text{点阵} + \text{结构基元}
\]

\begin{figure}[H]\centering
\includegraphics[width=0.9\textwidth,keepaspectratio]{点阵-结构基元-晶胞.关系图.png}
\caption{图 点阵·结构基元·晶胞三要素关系图——核心关系式的可视化表达：晶体结构 = 点阵 + 结构基元；晶胞是最小重复单元，平移复制得到整个晶体。左下角给出NaCl、CsCl、金刚石的实例对比}
\end{figure}

\textbf{点阵的基本特征}：点阵是按周期性规律在空间排布的一组无限多个点。按连接其中任意两点的向量进行平移时，能使点阵复原。点阵中每个点具有相同的周围环境。

\begin{figure}[H]\centering
\includegraphics[width=0.9\textwidth,keepaspectratio]{13-2c-crystal-structure-to-lattice-abstraction.png}
\caption{图 从晶体结构抽象为点阵的过程——将每个结构基元用一个几何点（点阵点）表示，无数点阵点按周期性排列即构成空间点阵，连接点阵点即得空间格子（晶格）（普化原理第4版 图13.2c）}
\end{figure}

\subsubsection{2.2 平面点阵与正当单位}\label{ux5e73ux9762ux70b9ux9635ux4e0eux6b63ux5f53ux5355ux4f4d}

平面点阵可划分为无数并置的平行四边形（平面格子）。划分规则：

{\def\LTcaptype{none} % do not increment counter
\begin{longtable}[]{@{}cc@{}}
\toprule\noalign{}
格子类型 & 每个平行四边形平均含点阵点数 \\
\midrule\noalign{}
\endhead
\bottomrule\noalign{}
\endlastfoot
素格子 & 1 个 \\
复格子 & \(>1\) 个 \\
\end{longtable}
}

\textbf{正当平面单位的选取标准}（按优先级）：
1. 与平面点阵对称性一致的平行四边形
2. 对称性尽可能高（直角尽可能多）
3. 包含点阵点数目尽可能少（面积尽可能小）

按此标准，平面格子共有 \textbf{四种形状、五种型式}：正方形、六方、矩形、带心矩形、平行四边形。其中仅带心矩形为复格子，其余均为素格子。

\begin{quote}
⚠️ \textbf{易错点}：「点阵点 = 原子位置？」------不是！点阵点是抽象的几何点，代表一个结构基元的位置。原子是实际物质。NaCl晶体中每个点阵点对应1个Na\textsuperscript{+}+1个Cl\textsuperscript{-}（1个NaCl），而不是1个原子。点阵点之间可能是原子、离子或分子，具体由结构基元决定。
\end{quote}

\subsubsection{2.3 空间点阵与14种布拉维格子}\label{ux7a7aux95f4ux70b9ux9635ux4e0e14ux79cdux5e03ux62c9ux7ef4ux683cux5b50}

空间点阵可划分为无数并置的平行六面体（空间格子）。平行六面体单位中：

{\def\LTcaptype{none} % do not increment counter
\begin{longtable}[]{@{}cc@{}}
\toprule\noalign{}
点阵点位置 & 对每个单位的平均贡献 \\
\midrule\noalign{}
\endhead
\bottomrule\noalign{}
\endlastfoot
顶角 & \(1/8\) \\
棱上 & \(1/4\) \\
面上 & \(1/2\) \\
内部 & \(1\) \\
\end{longtable}
}

\textbf{正当空间单位的选取标准}：
1. 符合相应空间点阵的对称性
2. 棱与棱之间的直角数目尽量多
3. 体积尽量小

法国晶体学家布拉维（Bravais）证明：晶体的点阵共有 \textbf{7 种形状、14 种型式}，即 \textbf{14 种布拉维格子}。

{\def\LTcaptype{none} % do not increment counter
\begin{longtable}[]{@{}ccccc@{}}
\toprule\noalign{}
晶系 & 简单（P） & 体心（I） & 面心（F） & 底心（C） \\
\midrule\noalign{}
\endhead
\bottomrule\noalign{}
\endlastfoot
三斜 & ✅ & --- & --- & --- \\
单斜 & ✅ & --- & --- & ✅ \\
正交 & ✅ & ✅ & ✅ & ✅ \\
四方 & ✅ & ✅ & --- & --- \\
六方 & ✅ & --- & --- & --- \\
菱方（三方） & ✅ & --- & --- & --- \\
立方 & ✅ & ✅ & ✅ & --- \\
\end{longtable}
}

\begin{figure}[H]\centering
\includegraphics[width=0.9\textwidth,keepaspectratio]{14种布拉维格子矩阵.png}
\caption{图 14种布拉维格子——晶系×带心方式记忆矩阵。绿色标记存在的组合，灰色标记不存在的组合。底部附记忆口诀和总数统计：1+2+4+2+1+1+3=14}
\end{figure}

\textbf{记忆要点}：并非每个晶系都有全部四种带心方式。例如立方晶系只有 P、I、F 三种（无底心）；四方晶系只有 P、I 两种（无底心和面心------因为四方面心可由更小的四方体心代替，四方底心可由更小的简单四方代替）。

\begin{figure}[H]\centering
\includegraphics[width=0.9\textwidth,keepaspectratio]{13-3c-bravais-lattice-types-cubic-hexagonal.png}
\caption{图 14 种布拉维格子中四种常见类型（从左至右：简单立方 P、体心立方 I、面心立方 F、简单六方）— 立方晶系中 P、I、F 三种型式均存在，六方晶系仅简单型式（普化原理第4版 图13.3）}
\end{figure}

\begin{center}\rule{0.5\linewidth}{0.5pt}\end{center}

\subsection{三、七大晶系}\label{ux4e09ux4e03ux5927ux6676ux7cfb}

\subsubsection{3.1 晶胞参数}\label{ux6676ux80deux53c2ux6570}

\begin{figure}[H]\centering
\includegraphics[width=0.9\textwidth,keepaspectratio]{七大晶系判定决策树.png}
\caption{图 七大晶系判定决策树——从最高次对称轴出发，通过逐级判断确定晶系归属：6次轴→六方；4次轴+4×3次轴→立方 vs 4次轴→四方；3次轴→三方；2次轴→正交(3个) vs 单斜(1个)；无→三斜。底部附七大晶系参数速查表和判断要点}
\end{figure}

晶胞的大小和形状由 \textbf{晶胞参数} 决定：三个棱长 \(a\)、\(b\)、\(c\)，以及棱间的夹角 \(\alpha\)（\(b\) 与 \(c\) 夹角）、\(\beta\)（\(a\) 与 \(c\) 夹角）、\(\gamma\)（\(a\) 与 \(b\) 夹角）。

\subsubsection{3.2 七大晶系的晶胞参数特征}\label{ux4e03ux5927ux6676ux7cfbux7684ux6676ux80deux53c2ux6570ux7279ux5f81}

{\def\LTcaptype{none} % do not increment counter
\begin{longtable}[]{@{}
  >{\centering\arraybackslash}p{(\linewidth - 8\tabcolsep) * \real{0.2273}}
  >{\raggedright\arraybackslash}p{(\linewidth - 8\tabcolsep) * \real{0.1818}}
  >{\raggedright\arraybackslash}p{(\linewidth - 8\tabcolsep) * \real{0.1818}}
  >{\centering\arraybackslash}p{(\linewidth - 8\tabcolsep) * \real{0.2273}}
  >{\raggedright\arraybackslash}p{(\linewidth - 8\tabcolsep) * \real{0.1818}}@{}}
\toprule\noalign{}
\begin{minipage}[b]{\linewidth}\centering
晶系
\end{minipage} & \begin{minipage}[b]{\linewidth}\raggedright
轴长关系
\end{minipage} & \begin{minipage}[b]{\linewidth}\raggedright
轴角关系
\end{minipage} & \begin{minipage}[b]{\linewidth}\centering
独立参数数
\end{minipage} & \begin{minipage}[b]{\linewidth}\raggedright
实例
\end{minipage} \\
\midrule\noalign{}
\endhead
\bottomrule\noalign{}
\endlastfoot
\textbf{立方} & \(a = b = c\) & \(\alpha = \beta = \gamma = 90^\circ\) & 1 & NaCl, CsCl, Cu, 金刚石 \\
\textbf{四方} & \(a = b \neq c\) & \(\alpha = \beta = \gamma = 90^\circ\) & 2 & TiO\textsubscript{2}（金红石）, CaC\textsubscript{2} \\
\textbf{六方} & \(a = b \neq c\) & \(\alpha = \beta = 90^\circ,\ \gamma = 120^\circ\) & 2 & Mg, Zn, 石墨 \\
\textbf{菱方（三方）} & \(a = b = c\) & \(\alpha = \beta = \gamma < 120^\circ,\ \neq 90^\circ\) & 2 & CaCO\textsubscript{3}（方解石） \\
\textbf{正交} & \(a \neq b \neq c\) & \(\alpha = \beta = \gamma = 90^\circ\) & 3 & 斜方硫 \\
\textbf{单斜} & \(a \neq b \neq c\) & \(\alpha = \gamma = 90^\circ,\ \beta \neq 90^\circ\) & 4 & 石膏 \\
\textbf{三斜} & \(a \neq b \neq c\) & \(\alpha \neq \beta \neq \gamma\) & 6 & 蓝晶石 \\
\end{longtable}
}

晶胞体积的通式：

\[
V = abc\sqrt{1 - \cos^2\alpha - \cos^2\beta - \cos^2\gamma + 2\cos\alpha\cos\beta\cos\gamma}
\]

对于立方晶系：\(V = a^3\)；对于正交晶系：\(V = abc\)。

\subsubsection{3.3 晶系的对称性判定}\label{ux6676ux7cfbux7684ux5bf9ux79f0ux6027ux5224ux5b9a}

在竞赛中，判断晶系优先依据 \textbf{对称性} 而非仅看轴长关系。最高对称轴的类型决定了晶系归属：

{\def\LTcaptype{none} % do not increment counter
\begin{longtable}[]{@{}cc@{}}
\toprule\noalign{}
最高对称轴 & 晶系 \\
\midrule\noalign{}
\endhead
\bottomrule\noalign{}
\endlastfoot
四个 \(C_3\) 轴 & 立方 \\
一个 \(C_4\) 或 \(\overline{4}\) 轴 & 四方 \\
一个 \(C_6\) 或 \(\overline{6}\) 轴 & 六方 \\
一个 \(C_3\) 或 \(\overline{3}\) 轴 & 三方（菱方） \\
三个互相垂直的 \(C_2\) 轴或镜面 & 正交 \\
一个 \(C_2\) 轴或镜面 & 单斜 \\
无（只有 \(C_1\) 或 \(i\)） & 三斜 \\
\end{longtable}
}

\begin{center}\rule{0.5\linewidth}{0.5pt}\end{center}

\subsection{四、晶胞}\label{ux56dbux6676ux80de}

\subsubsection{4.1 晶胞的基本特征}\label{ux6676ux80deux7684ux57faux672cux7279ux5f81}

\textbf{晶胞（unit cell）} 是晶体结构中的基本重复单位。整个晶体按晶胞在三维空间周期性地重复排列，即无数晶胞 \textbf{无隙并置} 堆砌而成。

\textbf{注意}：晶胞必须是平行六面体。从晶体中取六方棱柱作为晶胞是错误的------六方晶系的正当晶胞仍为平行六面体（\(a = b \neq c,\ \gamma = 120^\circ\)）。

\subsubsection{4.2 素晶胞与复晶胞}\label{ux7d20ux6676ux80deux4e0eux590dux6676ux80de}

{\def\LTcaptype{none} % do not increment counter
\begin{longtable}[]{@{}
  >{\centering\arraybackslash}p{(\linewidth - 8\tabcolsep) * \real{0.2174}}
  >{\centering\arraybackslash}p{(\linewidth - 8\tabcolsep) * \real{0.2174}}
  >{\raggedright\arraybackslash}p{(\linewidth - 8\tabcolsep) * \real{0.1739}}
  >{\centering\arraybackslash}p{(\linewidth - 8\tabcolsep) * \real{0.2174}}
  >{\raggedright\arraybackslash}p{(\linewidth - 8\tabcolsep) * \real{0.1739}}@{}}
\toprule\noalign{}
\begin{minipage}[b]{\linewidth}\centering
类型
\end{minipage} & \begin{minipage}[b]{\linewidth}\centering
符号
\end{minipage} & \begin{minipage}[b]{\linewidth}\raggedright
特征
\end{minipage} & \begin{minipage}[b]{\linewidth}\centering
每个晶胞含点阵点数
\end{minipage} & \begin{minipage}[b]{\linewidth}\raggedright
示例
\end{minipage} \\
\midrule\noalign{}
\endhead
\bottomrule\noalign{}
\endlastfoot
\textbf{素晶胞} & P & 晶体结构中最小的基本单元 & 1 & 简单立方（Po） \\
\textbf{体心晶胞} & I & 可作体心平移 \((1/2, 1/2, 1/2)\) & 2 & 金属钾（bcc） \\
\textbf{面心晶胞} & F & 可作面心平移 \((1/2, 1/2, 0)\) 等 & 4 & 金属铜（fcc） \\
\textbf{底心晶胞} & C & 可作底心平移 \((1/2, 1/2, 0)\) & 2 & 碘晶体（B底心） \\
\end{longtable}
}

\begin{figure}[H]\centering
\includegraphics[width=0.9\textwidth,keepaspectratio]{13-7a-simple-cubic-packing-layer-top-view.png}
\caption{图 三种基本晶胞类型对比（从左至右：简单立方堆积及其晶胞、体心立方晶胞、面心立方晶胞）— 简单立方 Po 型（配位数 6，空间利用率 52 \%），体心立方 bcc（配位数 8，68 \%），面心立方 fcc（配位数 12，74 \%）（普化原理第4版 图13.7、图13.10）}
\end{figure}

\textbf{体心平移检验}：将原晶胞框架的顶角平移到体心位置，若得到的新晶胞与原晶胞无差别，则为体心晶胞。

例如氯化铯（CsCl）虽然看起来有体心原子，但 Cs\textsuperscript{+} 与 Cl\textsuperscript{-} 环境不同，不能作体心平移 → \textbf{不是} 体心晶胞，而是简单立方晶胞。

\subsubsection{4.3 晶胞中原子的分数坐标}\label{ux6676ux80deux4e2dux539fux5b50ux7684ux5206ux6570ux5750ux6807}

晶胞中原子的位置用 \textbf{分数坐标（fractional coordinates）} 表示：\((x, y, z)\)，表示向量 \(x\boldsymbol{a} + y\boldsymbol{b} + z\boldsymbol{c}\)。

\begin{figure}[H]\centering
\includegraphics[width=0.9\textwidth,keepaspectratio]{分数坐标三维可视化.png}
\caption{图 分数坐标三维可视化——左：立方晶胞中标注原点(0,0,0)、对角点(1,1,1)、体心(½,½,½)等关键位置及坐标轴方向；右：常见位置坐标速查表与NaCl坐标示例。底部附写坐标的三个要点}
\end{figure}

\textbf{常见位置的分数坐标}：

{\def\LTcaptype{none} % do not increment counter
\begin{longtable}[]{@{}cc@{}}
\toprule\noalign{}
位置 & 坐标 \\
\midrule\noalign{}
\endhead
\bottomrule\noalign{}
\endlastfoot
顶角 & \((0, 0, 0)\) \\
体心 & \((1/2, 1/2, 1/2)\) \\
\(a\) 棱心 & \((1/2, 0, 0)\) \\
\(b\) 棱心 & \((0, 1/2, 0)\) \\
\(c\) 棱心 & \((0, 0, 1/2)\) \\
\(ab\) 面心 & \((1/2, 1/2, 0)\) \\
\(ac\) 面心 & \((1/2, 0, 1/2)\) \\
\(bc\) 面心 & \((0, 1/2, 1/2)\) \\
\end{longtable}
}

\textbf{注意}：分数坐标的取值区间为 \(0 \leqslant x, y, z < 1\)。由晶胞的平移性可知，坐标 \((1, 0, 0)\) 等价于 \((0, 0, 0)\)。

\subsubsection{4.4 晶胞中的原子计数}\label{ux6676ux80deux4e2dux7684ux539fux5b50ux8ba1ux6570}

\begin{figure}[H]\centering
\includegraphics[width=0.9\textwidth,keepaspectratio]{晶胞原子计数示意图.png}
\caption{图 晶胞原子计数——位置×贡献率。左：四种原子位置类型（顶点⅛、棱心¼、面心½、体心1）及共享说明；右：常见晶胞Z值计算（sc=1, bcc=2, fcc=4, hcp=2）。底部附易错提醒}
\end{figure}

{\def\LTcaptype{none} % do not increment counter
\begin{longtable}[]{@{}ccl@{}}
\toprule\noalign{}
原子位置 & 对晶胞的贡献 & 计数示例 \\
\midrule\noalign{}
\endhead
\bottomrule\noalign{}
\endlastfoot
顶角 & \(1/8\) & 立方晶胞：\(8 \times 1/8 = 1\) \\
棱心 & \(1/4\) & 立方晶胞：\(12 \times 1/4 = 3\) \\
面心 & \(1/2\) & 立方晶胞：\(6 \times 1/2 = 3\) \\
体内 & \(1\) & --- \\
\end{longtable}
}

\textbf{示例------金刚石晶胞}：包含 8 个碳原子（顶角 \(8 \times 1/8\) + 面心 \(6 \times 1/2\) + 体内 4 个），坐标分别为：
- \((0, 0, 0),\ (1/2, 1/2, 0),\ (1/2, 0, 1/2),\ (0, 1/2, 1/2)\)
- \((3/4, 1/4, 1/4),\ (1/4, 3/4, 1/4),\ (1/4, 1/4, 3/4),\ (3/4, 3/4, 3/4)\)

\begin{figure}[H]\centering
\includegraphics[width=0.9\textwidth,keepaspectratio]{13-15c-diamond-unit-cell-atom-count.png}
\caption{图 金刚石的晶体结构和面心立方晶胞（左：金刚石中 C 原子以 sp3 杂化形成的三维网状结构；右：从金刚石结构中取出的面心立方晶胞，含 8 个碳原子，其中 4 个在顶点/面心、4 个在体内四面体空隙）（普化原理第4版 图13.15）}
\end{figure}

\begin{center}\rule{0.5\linewidth}{0.5pt}\end{center}

\subsection{五、等同点系}\label{ux4e94ux7b49ux540cux70b9ux7cfb}

\subsubsection{5.1 定义}\label{ux5b9aux4e49}

晶体结构中，若一组点从其中各个点出发，沿同一指定方向观察到的物质环境和几何环境都相同，则这些点属于同一个 \textbf{等同点系（equivalent point system）}。

\subsubsection{5.2 拆分规则}\label{ux62c6ux5206ux89c4ux5219}

{\def\LTcaptype{none} % do not increment counter
\begin{longtable}[]{@{}cl@{}}
\toprule\noalign{}
步骤 & 操作 \\
\midrule\noalign{}
\endhead
\bottomrule\noalign{}
\endlastfoot
① & 先把不同的原子相互拆开（不同元素肯定不属于同一等同点系） \\
② & 相同的原子若周围环境不同，需拆分为多个等同点系 \\
③ & 拆分到各等同点系上原子的排列方式已相同时停止 \\
\end{longtable}
}

\subsubsection{5.3 应用示例}\label{ux5e94ux7528ux793aux4f8b}

\textbf{CsCl 晶体}：拆分为 2 个等同点系------Cl\textsuperscript{-} 按简单立方排列，Cs\textsuperscript{+} 按简单立方排列。每个等同点系上离子的排列方式即为点阵------简单立方点阵。从两个等同点系各取 1 个离子（1 个 Cl\textsuperscript{-} + 1 个 Cs\textsuperscript{+}）= \textbf{结构基元}。

\textbf{金刚石晶体}：有两类环境不同的碳原子------晶胞顶点/面心处的 CⅠ 和体内四面体空隙中的 CⅡ。拆分为 2 个等同点系，各自按面心立方排列。点阵为 \textbf{面心立方点阵}。结构基元 = 1 个 CⅠ + 1 个 CⅡ = \textbf{2 个碳原子}。

\begin{figure}[H]\centering
\includegraphics[width=0.9\textwidth,keepaspectratio]{13-11a-nacl-cscl-zns-three-ab-structures.png}
\caption{图 三种典型的 AB 型离子晶体结构（从左至右：NaCl 型、CsCl 型、立方 ZnS 型）— NaCl 中 Cl- 面心立方密堆积，Na+ 填全部八面体空隙；CsCl 中 Cl- 简单立方堆积，Cs+ 填立方体空隙；ZnS 中 S2- 面心立方密堆积，Zn2+ 填半数四面体空隙（普化原理第4版 图13.11）}
\end{figure}

\subsubsection{5.4 等同点系与点阵的区别}\label{ux7b49ux540cux70b9ux7cfbux4e0eux70b9ux9635ux7684ux533aux522b}

\textbf{易错点}：不能认为点阵是从等同点系抽象出来的。任何晶体都只能抽象出 \textbf{唯一确定} 的点阵，但可能有多个等同点系。一个点阵点代表的结构基元必须从每个等同点系上各取 1 个原子集合而成。

\begin{quote}
🧠 \textbf{认知冲突}：「等同点系与点阵」------这是晶体学中最易混淆的概念之一。\textbf{点阵}是抽象几何点，\textbf{等同点系}是晶体中实际原子位置的集合。点阵是''骨架''（几何位置关系），等同点系是''骨架上的内容''。不同等同点系可以有相同的排列方式（如CsCl中Cl\textsuperscript{-}和Cs\textsuperscript{+}都按简单立方排列），但它们属于不同的等同点系。
\end{quote}

\begin{center}\rule{0.5\linewidth}{0.5pt}\end{center}

\subsection{六、对称元素基础}\label{ux516dux5bf9ux79f0ux5143ux7d20ux57faux7840}

本节的对称元素概念与本讲义配套的 04-课件/学生讲义/2026-06-23-分子结构和性质\#五、对称性基础 中分子对称性的内容互补------分子对称性属于点对称性（点操作），晶体对称性还涉及平移对称性。

\subsubsection{6.1 晶体中的对称元素}\label{ux6676ux4f53ux4e2dux7684ux5bf9ux79f0ux5143ux7d20}

{\def\LTcaptype{none} % do not increment counter
\begin{longtable}[]{@{}
  >{\centering\arraybackslash}p{(\linewidth - 6\tabcolsep) * \real{0.2778}}
  >{\centering\arraybackslash}p{(\linewidth - 6\tabcolsep) * \real{0.2778}}
  >{\raggedright\arraybackslash}p{(\linewidth - 6\tabcolsep) * \real{0.2222}}
  >{\raggedright\arraybackslash}p{(\linewidth - 6\tabcolsep) * \real{0.2222}}@{}}
\toprule\noalign{}
\begin{minipage}[b]{\linewidth}\centering
对称元素
\end{minipage} & \begin{minipage}[b]{\linewidth}\centering
符号
\end{minipage} & \begin{minipage}[b]{\linewidth}\raggedright
操作
\end{minipage} & \begin{minipage}[b]{\linewidth}\raggedright
说明
\end{minipage} \\
\midrule\noalign{}
\endhead
\bottomrule\noalign{}
\endlastfoot
\textbf{旋转轴} & \(C_n\) & 绕轴旋转 \(360^\circ/n\) 复原 & \(n = 1, 2, 3, 4, 6\)（晶体中无 \(C_5\) 及 \(n > 6\) 的旋转轴） \\
\textbf{镜面} & \(\sigma\) & 通过镜面反映 & 镜面将晶体分为互为镜像的两部分 \\
\textbf{对称中心} & \(i\) & 通过中心反演 & 将 \((x, y, z)\) 变换为 \((-x, -y, -z)\) \\
\textbf{映轴} & \(S_n\) & 旋转 \(360^\circ/n\) 后对垂直面反映 & \(S_4\) 在晶体中出现（如 CH\textsubscript{4} 的正四面体对称性） \\
\textbf{反轴} & \(\overline{n}\) & 旋转 \(360^\circ/n\) 后反演 & \(\overline{1} = i\)，\(\overline{2} = \sigma\)（镜面） \\
\end{longtable}
}

\subsubsection{6.2 晶体的对称性限制}\label{ux6676ux4f53ux7684ux5bf9ux79f0ux6027ux9650ux5236}

\begin{figure}[H]\centering
\includegraphics[width=0.9\textwidth,keepaspectratio]{对称性限制定律.png}
\caption{图 晶体对称性限制定律——左：几何直观（五边形不能铺满平面）；右：数学证明（旋转矩阵迹条件）。底部总结三大影响：七大晶系对称性来源、准晶的发现（2011诺贝尔化学奖）、点阵对称性对结构基元排列的限制}
\end{figure}

晶体中旋转轴的轴次只能为 \(n = 1, 2, 3, 4, 6\)，不可能出现 \(C_5\) 及 \(n > 6\) 的旋转轴。这源于平移对称性的约束------称为 \textbf{晶体对称性限制定律}。

\begin{quote}
🧠 \textbf{认知冲突}：「准晶 (Quasicrystal) 打破了对称性限制定律？」------不完全是！1982年 Shechtman 在急冷 Al-Mn 合金中发现了5次对称衍射图案，这似乎''打破''了传统晶体学定律。但准晶并非周期性晶体------它是有序但非周期性的结构（利用彭罗斯铺砌的数学原理）。这一发现迫使国际晶体学联合会重新定义''晶体''，Shechtman 也因此获得2011年诺贝尔化学奖。
\end{quote}

\subsubsection{6.3 晶系与对称性的对应}\label{ux6676ux7cfbux4e0eux5bf9ux79f0ux6027ux7684ux5bf9ux5e94}

晶体的七个晶系由其最高对称轴唯一确定（见 §3.3 的表）。竞赛中常见的对应关系：

{\def\LTcaptype{none} % do not increment counter
\begin{longtable}[]{@{}lcl@{}}
\toprule\noalign{}
对称性特征 & 晶系 & 典型实例 \\
\midrule\noalign{}
\endhead
\bottomrule\noalign{}
\endlastfoot
有 4 个 \(C_3\) 轴 & 立方 & 金刚石、NaCl、Cu \\
有 1 个 \(C_4\) 轴 & 四方 & TiO\textsubscript{2}（金红石）、SnO\textsubscript{2} \\
有 1 个 \(C_6\) 轴 & 六方 & 石墨、Zn（hcp） \\
有 1 个 \(C_3\) 轴 & 三方 & CaCO\textsubscript{3}（方解石） \\
\end{longtable}
}

\textbf{深化}：晶体的完整对称性分类（32 种点群、230 种空间群）属于第二轮/第三轮结构化学深化内容，本讲义不展开。

\begin{center}\rule{0.5\linewidth}{0.5pt}\end{center}

\subsection{七、规律与方法}\label{ux4e03ux89c4ux5f8bux4e0eux65b9ux6cd5}

\subsubsection{7.1 晶体学分析的三步流程}\label{ux6676ux4f53ux5b66ux5206ux6790ux7684ux4e09ux6b65ux6d41ux7a0b}

\begin{verbatim}
宏观晶体外形 → 对称性判断 → 晶系归属
    ↓
抽象点阵 → 点阵型式（P/I/F/C）→ 布拉维格子
    ↓
确定晶胞 → 原子坐标 → 原子计数 → 结构基元
\end{verbatim}

\subsubsection{7.2 常见误区}\label{ux5e38ux89c1ux8befux533a}

\textbf{误区 1}：``晶胞就是最小的结构单元''------晶胞是点阵的基本重复单位，但不一定是最小的（素晶胞才是最小的）。复晶胞（I、F、C）是素晶胞的多倍体。
\textbf{误区 2}：``CsCl 是体心立方''------CsCl 看似体心立方，但 Cs\textsuperscript{+} 与 Cl\textsuperscript{-} 环境不同，不能作体心平移，应归为简单立方点阵。
\textbf{误区 3}：``有 \(C_5\) 轴的分子不能有晶体''------分子可以具有 \(C_5\) 轴（如二茂铁），但晶体中不允许出现 \(C_5\) 轴。具有 \(C_5\) 轴的分子在结晶时，其分子对称性与晶体的平移对称性不兼容。
\textbf{误区 4}：``所有晶系都有 14 种布拉维格子''------14 种是所有晶系的格子总和。每个晶系只有特定几种带心方式（如立方只有 P/I/F，四方只有 P/I）。

\begin{center}\rule{0.5\linewidth}{0.5pt}\end{center}

\subsection{八、典型例题}\label{ux516bux5178ux578bux4f8bux9898}

\subsubsection{例 1 {[}晶系判断 ⭐⭐}\label{ux4f8b-1-ux6676ux7cfbux5224ux65ad}

\textbf{题目}：某晶体的晶胞参数为 \(a = 0.520\ \mathrm{nm},\ b = 0.520\ \mathrm{nm},\ c = 0.856\ \mathrm{nm},\ \alpha = \beta = \gamma = 90^\circ\)。判断该晶体属于什么晶系，有多少个独立晶胞参数。

\textbf{完整解答}：
- \(a = b \neq c\)，\(\alpha = \beta = \gamma = 90^\circ\) → \textbf{四方晶系}
- 独立参数：\(a\)（或 \(b\)）和 \(c\) → \textbf{2 个}

\begin{center}\rule{0.5\linewidth}{0.5pt}\end{center}

\subsubsection{例 2 {[}晶胞原子计数 ⭐⭐}\label{ux4f8b-2-ux6676ux80deux539fux5b50ux8ba1ux6570}

\textbf{题目}：金属钠的晶体结构为体心立方（bcc），每个晶胞含多少个钠原子？写出体心立方晶胞中所有原子的分数坐标。

\textbf{完整解答}：
- bcc 晶胞：顶角 \(8 \times 1/8 = 1\)，体心 \(1 \times 1 = 1\) → 共 \textbf{2 个} 原子
- 分数坐标：\((0, 0, 0)\) 和 \((1/2, 1/2, 1/2)\)

\begin{center}\rule{0.5\linewidth}{0.5pt}\end{center}

\subsubsection{例 3 {[}点阵型式判断 ⭐⭐⭐}\label{ux4f8b-3-ux70b9ux9635ux578bux5f0fux5224ux65ad}

\textbf{题目}：CaC\textsubscript{2} 晶体如下图所示（略），其中 C---C 键长为 0.120 nm，C\textsubscript{2}\textsuperscript{2-} 离子呈哑铃型且键轴平行于 c 轴排列。问 CaC\textsubscript{2} 的正当晶胞是否仍为面心立方？属于什么点阵型式？

\textbf{思路分析}：
1. CaC\textsubscript{2} 初看类似 NaCl 型结构，Na\textsuperscript{+} 替换为 Ca\textsuperscript{2+}，Cl\textsuperscript{-} 替换为 C\textsubscript{2}\textsuperscript{2-}
2. 但 C\textsubscript{2}\textsuperscript{2-} 为哑铃型，其键轴平行于 c 轴 → 破坏了立方对称性
3. 需要从对称性角度判断晶系归属和点阵型式

\textbf{完整解答}：
- C\textsubscript{2}\textsuperscript{2-} 键轴平行于 c 轴，使 c 轴膨胀，\(a \neq c\)；对称性从立方降低为四方
- 立方晶系的四个 \(C_3\) 轴被破坏，但 c 轴上的 \(C_4\) 轴保留 → \textbf{四方晶系}
- 面心四方可由更小的体心四方代替 → 正当晶胞为 \textbf{体心四方（I）}
- 所以 CaC\textsubscript{2} 的正当晶胞 \textbf{不是} 面心立方，而是体心四方

\textbf{答案}：不是面心立方；四方体心（I）点阵。

\begin{center}\rule{0.5\linewidth}{0.5pt}\end{center}

\subsubsection{例 4 {[}等同点系 ⭐⭐⭐}\label{ux4f8b-4-ux7b49ux540cux70b9ux7cfb}

\textbf{题目}：试分析 Cu\textsubscript{2}O（赤铜矿）的晶胞是否为体心立方？已知 Cu\textsubscript{2}O 晶胞中 O 原子坐标为 \((0, 0, 0)\) 和 \((1/2, 1/2, 1/2)\)；Cu 原子坐标为 \((3/4, 1/4, 1/4)\)、\((1/4, 3/4, 1/4)\)、\((1/4, 1/4, 3/4)\)、\((3/4, 3/4, 3/4)\)。

\textbf{完整解答}：
- 进行体心平移：将各原子坐标 \(+(1/2, 1/2, 1/2)\)
- O 原子新坐标：\((1/2, 1/2, 1/2)\) 和 \((0, 0, 0)\)------与原 O 原子坐标相同
- Cu 原子新坐标：\((1/4, 3/4, 3/4)\)、\((3/4, 1/4, 3/4)\)、\((3/4, 3/4, 1/4)\)、\((1/4, 1/4, 1/4)\)------与原 Cu 原子坐标不同
- 体心平移后 Cu 原子位置发生了变化 → \textbf{不具备体心平移特征}
- \textbf{结论}：Cu\textsubscript{2}O 的该晶胞 \textbf{不是} 体心立方晶胞

\textbf{启示}：判断是否为体心晶胞，不能仅看''体心是否有原子''，核心检验标准是 \textbf{体心平移后晶体结构是否不变}。

\begin{center}\rule{0.5\linewidth}{0.5pt}\end{center}

\subsubsection{例 5 {[}对称性判断结构的应用 ⭐⭐⭐}\label{ux4f8b-5-ux5bf9ux79f0ux6027ux5224ux65adux7ed3ux6784ux7684ux5e94ux7528}

\textbf{题目}：某晶体的粉末衍射数据显示其属于立方晶系，晶胞参数 \(a = 543\ \mathrm{pm}\)。已知该晶体中原子排列的对称性显示其具有面心平移特征（F），判读其每个晶胞含 4 个原子。若该晶体的密度为 \(2.33\ \mathrm{g\cdot cm^{-3}}\)，求其摩尔质量，并推断可能的物质。

\textbf{完整解答}：
- 面心立方晶胞含 \(8 \times 1/8 + 6 \times 1/2 = 4\) 个原子，与题述吻合
- 密度公式：\(\rho = \dfrac{nM}{N_A a^3}\)
- \(M = \dfrac{\rho N_A a^3}{n} = \dfrac{2.33 \times 6.022\times 10^{23} \times (5.43\times 10^{-8})^3}{4}\)
- \(M \approx 28.1\ \mathrm{g\cdot mol^{-1}}\) → \textbf{硅（Si）}

\begin{center}\rule{0.5\linewidth}{0.5pt}\end{center}

\subsection{九、拓展与联系}\label{ux4e5dux62d3ux5c55ux4e0eux8054ux7cfb}

\subsubsection{9.1 与竞赛真题的对接}\label{ux4e0eux7adeux8d5bux771fux9898ux7684ux5bf9ux63a5}

{\def\LTcaptype{none} % do not increment counter
\begin{longtable}[]{@{}lcl@{}}
\toprule\noalign{}
考点 & 频率 & 说明 \\
\midrule\noalign{}
\endhead
\bottomrule\noalign{}
\endlastfoot
晶系判断 & 每年 1 题 & 根据晶胞参数或对称性判断晶系 \\
晶胞原子计数 & 每年 1-2 题 & 顶角 \(1/8\)、面心 \(1/2\)、体内 \(1\) \\
分数坐标 & 隔年 1 题 & 常与密度计算、化学式推断结合 \\
点阵型式判断 & 隔年 1 题 & 需要通过等同点系或平移检验判断 \\
对称元素与晶系 & 隔年 1 题 & 由对称元素反推晶系归属 \\
\end{longtable}
}

\subsubsection{9.2 跨专题联系}\label{ux8de8ux4e13ux9898ux8054ux7cfb}

\begin{itemize}
\tightlist
\item
  04-课件/学生讲义/2026-06-23-分子结构和性质\#五、对称性基础 → 分子对称性基础是晶体对称性的前置
\item
  04-课件/学生讲义/2026-06-23-金属键与金属晶体 → 金属密堆积（ccp/hcp/bcc）是晶胞概念的应用
\item
  晶体结构 → 晶体学 KP 详解
\item
  晶胞 → 晶胞参数与分数坐标
\end{itemize}

\subsubsection{9.3 深化方向（后续轮次）}\label{ux6df1ux5316ux65b9ux5411ux540eux7eedux8f6eux6b21}

\begin{itemize}
\tightlist
\item
  \textbf{X 射线衍射与 Bragg 方程}------第二轮晶体结构专题，需要波动光学基础
\item
  \textbf{空间群完整分类（230 种）与点群（32 种）}------第三轮结构化学
\item
  \textbf{群论在晶体对称性分析中的应用}------决赛/结构化学深化
\item
  \textbf{复杂结构（钙钛矿/尖晶石/反尖晶石）}------第二轮无机结构深化
\end{itemize}

\begin{center}\rule{0.5\linewidth}{0.5pt}\end{center}

\subsection{十、本节总结 · 速查}\label{ux5341ux672cux8282ux603bux7ed3-ux901fux67e5}

{\def\LTcaptype{none} % do not increment counter
\begin{longtable}[]{@{}
  >{\raggedright\arraybackslash}p{(\linewidth - 4\tabcolsep) * \real{0.3333}}
  >{\raggedright\arraybackslash}p{(\linewidth - 4\tabcolsep) * \real{0.3333}}
  >{\raggedright\arraybackslash}p{(\linewidth - 4\tabcolsep) * \real{0.3333}}@{}}
\toprule\noalign{}
\begin{minipage}[b]{\linewidth}\raggedright
核心概念
\end{minipage} & \begin{minipage}[b]{\linewidth}\raggedright
关键规律
\end{minipage} & \begin{minipage}[b]{\linewidth}\raggedright
常见陷阱
\end{minipage} \\
\midrule\noalign{}
\endhead
\bottomrule\noalign{}
\endlastfoot
\textbf{晶体定义} & 长程有序 + 周期性排列 & 非晶体也有短程有序但长程无序 \\
\textbf{晶体性质} & 自范性、各向异性、对称性、均一性 & 各向异性 ≠ 不均一 \\
\textbf{点阵} & 晶体结构 = 点阵 + 结构基元 & 等同点系 ≠ 点阵（晶体只有一个点阵） \\
\textbf{14 种布拉维格子} & 7 种晶系 × 带心方式（P/I/F/C） & 并非每种晶系都有全部 4 种 \\
\textbf{晶胞计数} & 顶角 \(1/8\)、棱心 \(1/4\)、面心 \(1/2\)、体内 \(1\) & CsCl 不是体心（体心平移条件不满足） \\
\textbf{七大晶系} & 由最高对称轴决定 & 仅凭轴长轴角判断有时不可靠 \\
\textbf{对称元素} & \(C_n\)、\(\sigma\)、\(i\)、\(S_n\)、\(\overline{n}\) & 晶体中无 \(C_5\)（对称性限制定律） \\
\textbf{素/复晶胞} & P = 1 个点阵点，I = 2，F = 4，C = 2 & 分数坐标 \((1, 0, 0)\) 等价于 \((0, 0, 0)\) \\
\end{longtable}
}

\begin{center}\rule{0.5\linewidth}{0.5pt}\end{center}

\subsection{十一、练习题}\label{ux5341ux4e00ux7ec3ux4e60ux9898}

\subsubsection{11.1 基础巩固}\label{ux57faux7840ux5de9ux56fa}

\begin{enumerate}
\def\labelenumi{\arabic{enumi}.}
\item
  \textbf{晶体性质}：请解释为什么晶体具有固定的熔点而非晶体没有。
\item
  \textbf{七大晶系判断}：根据以下晶胞参数判断所属晶系：

  \begin{enumerate}
  \def\labelenumii{(\arabic{enumii})}
  \tightlist
  \item
    \(a = b = c = 0.400\ \mathrm{nm},\ \alpha = \beta = \gamma = 90^\circ\)
  \item
    \(a = b = 0.500\ \mathrm{nm},\ c = 0.800\ \mathrm{nm},\ \alpha = \beta = 90^\circ,\ \gamma = 120^\circ\)
  \item
    \(a = 0.700\ \mathrm{nm},\ b = 0.700\ \mathrm{nm},\ c = 0.900\ \mathrm{nm},\ \alpha = \beta = \gamma = 90^\circ\)
  \end{enumerate}
\item
  \textbf{晶胞计数}：计算面心立方（fcc）晶胞中有多少个原子。写出 fcc 晶胞中所有原子的分数坐标。
\end{enumerate}

\subsubsection{11.2 提高训练}\label{ux63d0ux9ad8ux8badux7ec3}

\begin{enumerate}
\def\labelenumi{\arabic{enumi}.}
\setcounter{enumi}{3}
\item
  \textbf{点阵型式判断}：某化合物的晶胞中，A 原子位于顶点，B 原子位于面心，C 原子位于体心。判断：

  \begin{enumerate}
  \def\labelenumii{(\arabic{enumii})}
  \tightlist
  \item
    该晶胞属于简单、体心、面心还是底心？
  \item
    写出该晶体的化学式。
  \item
    写出所有原子的分数坐标。
  \end{enumerate}
\item
  \textbf{等同点系应用}：已知金刚石晶胞中碳原子坐标为 \((0, 0, 0)\)、\((1/2, 1/2, 0)\)、\((1/2, 0, 1/2)\)、\((0, 1/2, 1/2)\)、\((3/4, 1/4, 1/4)\)、\((1/4, 3/4, 1/4)\)、\((1/4, 1/4, 3/4)\)、\((3/4, 3/4, 3/4)\)。试证明金刚石晶体的点阵为面心立方，结构基元含 2 个碳原子。
\item
  \textbf{对称性与晶系}：某晶体的晶胞有三个互相垂直的棱，长度分别为 \(a = 0.510\ \mathrm{nm}\)、\(b = 0.430\ \mathrm{nm}\)、\(c = 0.640\ \mathrm{nm}\)，棱间夹角均为 \(90^\circ\)。该晶体有几个对称元素？属于什么晶系？
\end{enumerate}

\subsubsection{11.3 挑战题}\label{ux6311ux6218ux9898}

\begin{enumerate}
\def\labelenumi{\arabic{enumi}.}
\setcounter{enumi}{6}
\tightlist
\item
  \textbf{综合推理}：金属钨为体心立方（bcc）结构，晶胞参数 \(a = 0.316\ \mathrm{nm}\)，密度 \(\rho = 19.3\ \mathrm{g\cdot cm^{-3}}\)。

  \begin{enumerate}
  \def\labelenumii{(\arabic{enumii})}
  \tightlist
  \item
    计算每个晶胞中的钨原子数。
  \item
    计算阿伏加德罗常数 \(N_A\)。
  \item
    计算空间利用率（堆积效率）。
  \end{enumerate}
\item
  \textbf{晶系归属判定}：已知 CaC\textsubscript{2} 的晶体结构在室温下为体心四方（I），晶胞参数 \(a = 0.387\ \mathrm{nm}\)、\(c = 0.637\ \mathrm{nm}\)。当温度升高到 \(450^\circ\mathrm{C}\) 以上时，CaC\textsubscript{2} 转变为面心立方结构，晶胞参数 \(a = 0.580\ \mathrm{nm}\)。试解释：

  \begin{enumerate}
  \def\labelenumii{(\arabic{enumii})}
  \tightlist
  \item
    为什么室温时 CaC\textsubscript{2} 为四方晶系而非立方晶系？
  \item
    高温时对称性升高为立方晶系，C\textsubscript{2}\textsuperscript{2-} 离子取向可能发生什么变化？
  \item
    计算两种晶型下 CaC\textsubscript{2} 的理论密度，比较大小。
  \end{enumerate}
\end{enumerate}

\begin{center}\rule{0.5\linewidth}{0.5pt}\end{center}

\subsection{十二、参考答案与提示}\label{ux5341ux4e8cux53c2ux8003ux7b54ux6848ux4e0eux63d0ux793a}

\textbf{11.1 基础巩固}

\begin{enumerate}
\def\labelenumi{\arabic{enumi}.}
\item
  晶体具有长程有序的周期性结构，所有原子/离子处于相同类型的环境中，加热时各处的原子同时获得足够能量脱离晶格位置 → 在固定温度全部熔化。非晶体无长程有序，各处结合能不同 → 随温度升高逐渐软化，无固定熔点。
\item
  \begin{enumerate}
  \def\labelenumii{(\arabic{enumii})}
  \tightlist
  \item
    立方晶系；(2) 六方晶系；(3) 四方晶系。
  \end{enumerate}
\item
  fcc 晶胞含 \(8 \times 1/8 + 6 \times 1/2 = 4\) 个原子。分数坐标：\((0, 0, 0)\)、\((1/2, 1/2, 0)\)、\((1/2, 0, 1/2)\)、\((0, 1/2, 1/2)\)。
\end{enumerate}

\textbf{11.2 提高训练}

\begin{enumerate}
\def\labelenumi{\arabic{enumi}.}
\setcounter{enumi}{3}
\item
  \textbf{提示}：(1) A 在顶点、B 在面心、C 在体心 → 顶点平移至体心后结构不变（A 仍到 A，B 仍到 B）→ 体心晶胞（I）。(2) 化学式：\(A_{8\times 1/8} B_{6\times 1/2} C_{1\times 1} = \mathrm{AB_3C}\)。(3) A: \((0,0,0)\)；B: \((1/2,1/2,0)\)、\((1/2,0,1/2)\)、\((0,1/2,1/2)\)；C: \((1/2,1/2,1/2)\)。
\item
  \textbf{提示}：前 4 个原子（CⅠ）坐标位于顶点和面心 → 面心立方排列。后 4 个原子（CⅡ）坐标相对于 CⅠ偏移 \((1/4, 1/4, 1/4)\) → 亦为面心立方排列。两类碳原子环境不同，分属 2 个等同点系 → 点阵为面心立方。每个面心立方点阵点对应 1 个 CⅠ + 1 个 CⅡ = 2 个 C 原子 → 结构基元含 2 个碳原子。
\item
  \textbf{提示}：\(a \neq b \neq c\)，\(\alpha = \beta = \gamma = 90^\circ\) → 正交晶系。正交晶系在三个互相垂直的方向上各有 \(C_2\) 轴（或镜面）。
\end{enumerate}

\textbf{11.3 挑战题}

\begin{enumerate}
\def\labelenumi{\arabic{enumi}.}
\setcounter{enumi}{6}
\item
  \textbf{提示}：(1) bcc 晶胞含 2 个钨原子。(2) 由 \(\rho = \dfrac{2M}{N_A a^3}\)，代入 \(M_{\mathrm{W}} = 183.84\ \mathrm{g\cdot mol^{-1}}\)、\(a = 0.316\ \mathrm{nm} = 3.16\times 10^{-8}\ \mathrm{cm}\)，可解得 \(N_A\)。(3) bcc 空间利用率 = 68\%。
\item
  \textbf{提示}：(1) 室温时 C\textsubscript{2}\textsuperscript{2-} 的 C---C 键轴平行于 c 轴排列 → 破坏了立方晶系的四个 \(C_3\) 轴对称性，保留 \(C_4\) 轴 → 四方晶系。(2) 高温时 C\textsubscript{2}\textsuperscript{2-} 热运动加剧，C---C 键轴取向无序化（在各方向上随机分布）→ 宏观上呈现出立方对称性。(3) 计算密度 \(\rho = nM/N_A V\)，注意两种晶型晶胞所含 CaC\textsubscript{2} 的式量数不同（面心立方时 \(Z = 4\)，体心四方时 \(Z = 2\)），对比结果可知密度变化。
\end{enumerate}

\begin{center}\rule{0.5\linewidth}{0.5pt}\end{center}

\subsection{十三、延伸阅读}\label{ux5341ux4e09ux5ef6ux4f38ux9605ux8bfb}

\begin{itemize}
\tightlist
\item
  晶体结构 --- 晶体学基础 KP
\item
  晶胞 --- 晶胞参数与分数坐标详解
\item
  布拉维格子 --- 14 种格子的系统介绍
\item
  04-课件/学生讲义/2026-06-23-金属键与金属晶体 --- 金属密堆积的几何计算
\item
  04-课件/学生讲义/2026-06-23-分子结构和性质\#五、对称性基础 --- 分子对称性入门
\item
  专题-晶体结构计算 --- 专题知识系统总结
\item
  上海中学竞赛课程·化学·第一分册·第六讲
\end{itemize}

\begin{center}\rule{0.5\linewidth}{0.5pt}\end{center}

\emph{对应第一轮专题 12：专题-晶体结构计算（备课框架 v2.2 §十二）。}

\end{document}
